% Préambule du livre. Même mise en page que le livre de learn-lean
% (book/preambule.tex) : format 17 × 24 cm, texte en serif et titres en
% sans-serif, encadrés à icônes pour les notes et les mises en garde.
%
% Le document se compile avec XeTeX (donc avec tectonic) : fontspec et
% awesomebox l'exigent.
\documentclass[11pt,DIV=12,BCOR=6mm,open=any,numbers=noenddot,chapterprefix=false]{scrbook}

% Format 17 × 24 cm, celui des livres scientifiques français.
\usepackage[paperwidth=17cm,paperheight=24cm,
  inner=2.1cm,outer=1.9cm,top=2.2cm,bottom=2.4cm]{geometry}

% Le texte en XCharter et les titres en Fira Sans. L'ordre compte : les
% paquets de polices avant fontspec.
\usepackage{XCharter}
\usepackage[scaled=0.94]{FiraSans}
\usepackage{fontspec}
% La chasse fixe porte les commandes et les noms de fichiers.
\setmonofont{DejaVuSansMono.ttf}[Scale=0.82]
\usepackage[french]{babel}
\usepackage{amsmath,amssymb,amsthm}
\usepackage{microtype}

% Les encadrés à icônes : note, astuce, mise en garde, avertissement.
\usepackage{awesomebox}
\usepackage[most]{tcolorbox}

\usepackage{tikz}
\usetikzlibrary{calc,positioning,arrows.meta}

\usepackage[hidelinks]{hyperref}

% Les noms de fichiers et de commandes ne se coupent pas : mieux vaut des
% espacements plus lâches que des lignes qui débordent dans la marge.
\emergencystretch=2em

\DeclareTOCStyleEntry[numwidth=2.4em]{tocline}{section}
\DeclareTOCStyleEntry[numwidth=3.6em]{tocline}{subsection}

% Titres et en-têtes en sans-serif, comme dans un manuel technique.
\setkomafont{disposition}{\sffamily}
\setkomafont{descriptionlabel}{\sffamily\bfseries}
\addtokomafont{caption}{\small}
\usepackage[automark]{scrlayer-scrpage}
\pagestyle{scrheadings}
\clearpairofpagestyles
\automark[section]{chapter}
\ihead*{\small\sffamily\headmark}
\ohead*{\small\sffamily\pagemark}

% Les exemples, composés comme dans le livre de learn-lean : bloc détaché,
% filet et fond tramés.
\theoremstyle{definition}
\newtheorem{exemple}{Exemple}[chapter]
\newtheorem{regle}[exemple]{Règle}
\definecolor{tonexemple}{HTML}{B45309}
\tcolorboxenvironment{exemple}{
  enhanced, breakable, sharp corners, boxrule=0pt,
  colback=tonexemple!4, borderline west={2.5pt}{0pt}{tonexemple!70},
  left=10pt, right=8pt, top=6pt, bottom=6pt,
  before skip=10pt, after skip=10pt,
}
\definecolor{tonregle}{HTML}{1D4ED8}
\tcolorboxenvironment{regle}{
  enhanced, breakable, sharp corners, boxrule=0pt,
  colback=tonregle!4, borderline west={2.5pt}{0pt}{tonregle!70},
  left=10pt, right=8pt, top=6pt, bottom=6pt,
  before skip=10pt, after skip=10pt,
}

% Ce que ce livre ajoute au préambule de learn-lean : des commandes de
% terminal, composées en chasse fixe sur fond gris, sans numérotation.
\newtcblisting{terminal}{
  listing only, breakable, sharp corners, boxrule=0pt,
  colback=black!4, borderline west={2.5pt}{0pt}{black!40},
  left=8pt, right=6pt, top=4pt, bottom=4pt,
  before skip=8pt, after skip=8pt,
  listing options={basicstyle=\ttfamily\small, columns=fullflexible,
    keepspaces=true, breaklines=true, showstringspaces=false},
}
\newcommand{\fichier}[1]{\texttt{#1}}
\newcommand{\cmd}[1]{\texttt{#1}}

\title{%
  Mathématiques outillées\\
  \large le génie logiciel au service de la recherche}
\author{Michel Hua \\[8pt]
  \normalsize\normalfont avec Claude (Anthropic) comme assistant}
\date{}
